%
% Chapter 6: Technology Comparison
%
\chapter{Technology Comparison}
\label{chap:tech-comparison}
\setlength{\parskip}{1em}

This chapter explains the main technology choices made during Snowdrop’s development. Each decision is evaluated against three constraints that shaped the project from the beginning: predictable tail latency, feasibility within a proof-of-concept scope, and compatibility with the PostgreSQL ecosystem.

\section{Orchestrator Language}

\begin{table}[H]
\centering
\caption{Orchestrator language comparison}
\label{tab:lang-comparison}
\renewcommand{\arraystretch}{1.4}
\begin{tabular}{|l|p{4cm}|p{5.5cm}|}
\hline
\textbf{Option} & \textbf{Key property} & \textbf{Verdict} \\
\hline
\textbf{Rust} & Zero GC, Tokio async runtime, full crate ecosystem & \textbf{Chosen} \\
\hline
Go & Simple concurrency, used by CockroachDB & GC produces 1--5\,ms pauses - up to 17\% of the 30\,ms p99 budget \\
\hline
Zig & Manual memory, zero-cost C interop & No production async runtime, no pgwire/Bolt/RESP crates \\
\hline
\end{tabular}
\end{table}

Rust was selected because it removes garbage collection from the critical path while still offering a mature async runtime through Tokio~\cite{tokio2025}. It also provides a production-ready crate ecosystem that directly matches Snowdrop’s needs: \texttt{pgwire} for PostgreSQL protocol support~\cite{pgwire2025}, \texttt{sqlparser} for SQL parsing~\cite{sqlparser2025}, \texttt{neo4rs} for Neo4j Bolt connectivity, and \texttt{bincode} for WAL serialization.

Go was considered because of its simplicity and its proven use in systems such as CockroachDB, but its garbage collector introduces pause times that are difficult to ignore when the target is sub-30\,ms p99 latency. Zig, on the other hand, offers low-level control and efficient C interoperability, but at the time of implementation it lacked both a mature async runtime and the ecosystem support needed for PostgreSQL, Bolt, and RESP integration. In practice, Rust provided the strongest balance between performance requirements and implementation feasibility.

\section{Data Storage}

\begin{table}[H]
\centering
\caption{Data storage comparison}
\label{tab:storage-comparison}
\renewcommand{\arraystretch}{1.4}
\begin{tabular}{|l|p{3.5cm}|p{5.5cm}|}
\hline
\textbf{Option} & \textbf{Key property} & \textbf{Verdict} \\
\hline
\textbf{Embedded RocksDB} & In-process, zero network overhead & \textbf{Chosen} - eliminated 70\% of query latency \\
\hline
Apache Kvrocks & RESP-compatible, RocksDB on NVMe & Replaced - TCP round-trips dominated read path~\cite{kvrocks2025} \\
\hline
Vanilla Redis & Fastest in-memory KV store & All data in RAM - unviable at scale \\
\hline
\end{tabular}
\end{table}

This turned out to be the most consequential design decision in the project. The original version of Snowdrop used Apache Kvrocks as an external storage layer, which made early integration straightforward because of its Redis-compatible interface. However, once the system was profiled under analytical workloads, it became clear that the network boundary itself was the main bottleneck. Replacing Kvrocks with embedded RocksDB~\cite{matsunobu2020} reduced TPC-H Q6 latency from 165\,ms to 17\,ms~\cite{luo2020}.

The improvement came from removing TCP round-trips between the orchestrator and the storage engine. When read latency is expected to stay below a few dozen milliseconds, inter-process communication can consume more time than the actual query logic. For this reason, embedded RocksDB was adopted as the primary storage backend. At the same time, the earlier Kvrocks implementation was preserved behind the \texttt{StorageBackend} trait and can still be enabled through the \texttt{STORAGE\_BACKEND=kvrocks} environment variable for direct architectural comparison.

\section{Schema Authority and Secondary Index Placement}

\begin{table}[H]
\centering
\caption{Schema and secondary-index alternatives}
\label{tab:index-comparison}
\renewcommand{\arraystretch}{1.4}
\begin{tabular}{|l|p{3.5cm}|p{5.5cm}|}
\hline
\textbf{Approach} & \textbf{Property} & \textbf{Verdict} \\
\hline
\textbf{Neo4j for schema metadata + RocksDB for secondary indexes} & Clear separation of metadata and data paths & \textbf{Chosen} \\
\hline
B-tree-style index keys in RocksDB & In-process lookup near the row store & Adopted in the current implementation \\
\hline
Array property on Column node & Simple implementation & O(n) deserialization and poor scaling \\
\hline
\end{tabular}
\end{table}

Neo4j was chosen as the authority for schema metadata because its graph model fits Snowdrop’s needs in the schema layer, where relationships such as Database $\rightarrow$ Table $\rightarrow$ Column, foreign-key links, and constraint metadata are naturally represented as connected nodes~\cite{robinson2015, neo4j2025}. However, the current implementation no longer stores indexed values in Neo4j.

The earlier prototype explored representing secondary indexes as Neo4j graph objects. Profiling and later implementation work showed that the more practical design is to keep indexed values in RocksDB next to the durable row store. This eliminates Bolt/JVM traversal from the hot read path, simplifies WAL-coordinated writes, and keeps lookup semantics close to the storage engine that already owns row payloads.

Under the current design, Neo4j stores no row payloads and no per-value index entries. Its growth therefore depends on schema complexity rather than on indexed row volume, while RocksDB absorbs both the underlying data and the secondary-index key space.

\section{Client Protocol and SQL Parsing}

Snowdrop implements the PostgreSQL wire protocol through the \texttt{pgwire} crate, which allows standard PostgreSQL-compatible tools to connect without modification. This includes clients such as \texttt{psql}, JDBC, \texttt{asyncpg}, SQLAlchemy, Prisma, and others~\cite{postgresql2026}. This decision significantly reduced the adoption barrier for testing and experimentation. A custom protocol would have required separate client support for each language or framework before the system could even be evaluated properly.

SQL parsing is handled by \texttt{sqlparser-rs}, which is already used in production-oriented projects such as Apache DataFusion, Materialize, and GlueSQL. In Snowdrop, the parser is used only to produce an abstract syntax tree. Planning, execution, optimization decisions, and result serialization remain fully custom. This keeps the parser as a reusable front end while preserving architectural control over the rest of the database pipeline.

\section{Analytical Acceleration}

The columnar cache is the main feature that differentiates Snowdrop from more conventional storage engines~\cite{abadi2013}. Instead of maintaining data permanently in a column-oriented layout, Snowdrop extracts typed column arrays from RocksDB when a scan is first needed and then keeps those arrays in memory for later reuse. This makes it possible to accelerate analytical queries without fully abandoning row-oriented durable storage.

Predicate evaluation is performed with SIMD-accelerated comparison loops that operate over whole columns rather than one tuple at a time. The results are represented as \texttt{u64} bitmasks, which can then be combined using bitwise operations such as AND and OR~\cite{zhou2002}. This approach is especially effective for scan-heavy workloads with simple numeric predicates.

In practice, the columnar path delivers 0.8\,ms warm-cache latency on TPC-H Q6, compared to 1.5\,ms for PostgreSQL~16, which corresponds to a 1.9x improvement. Cold-cache execution is slower because the first scan must also populate the cache, adding a one-time cost of 8.2\,ms. After that, repeated scans benefit from the prebuilt in-memory representation. Query patterns that are not yet supported by the SIMD path fall back to relation or row-oriented evaluation.

\section{Summary}

Table~\ref{tab:tech-stack} summarizes the final technology stack used in Snowdrop.

\begin{table}[H]
\centering
\caption{Snowdrop technology stack}
\label{tab:tech-stack}
\renewcommand{\arraystretch}{1.4}
\small
\begin{tabular}{|l|l|l|}
\hline
\textbf{Component} & \textbf{Technology} & \textbf{Role in Snowdrop} \\
\hline
Language & Rust + Tokio & Orchestrator, async runtime \\
\hline
SQL parsing & sqlparser-rs & SQL string $\rightarrow$ AST \\
\hline
Wire protocol & pgwire crate & PostgreSQL client compatibility \\
\hline
Row storage & Embedded RocksDB & Durable key-value data \\
\hline
Schema metadata & Neo4j 5 & Graph-based metadata and introspection \\
\hline
Secondary indexes & Embedded RocksDB & In-process indexed lookup next to row storage \\
\hline
Analytical cache & Columnar + SIMD & Accelerated predicate evaluation \\
\hline
WAL format & bincode & Fast binary serialization \\
\hline
Containerization & Docker Compose & Development environment \\
\hline
Benchmark & TPC-H (scale factor 1) & Validation framework \\
\hline
\end{tabular}
\end{table}

Taken together, these choices reflect the broader logic of the project: use established components where they fit naturally, but combine them in a way that serves Snowdrop’s specific goals. Rather than trying to build every layer from scratch, the system relies on a focused set of technologies that support low-latency execution, practical implementation, and compatibility with familiar PostgreSQL tools.
